\chapter{Rigorous Mathematical Frameworks for the Open Problems}
\label{app:hard_analysis_frameworks}

\begin{tcolorbox}[colback=blue!5!white, colframe=blue!75!black, title=\textbf{Purpose of this Appendix}]
This appendix introduces the five specific mathematical frameworks employed to resolve the key open problems identified in the Ledger (Appendix~\ref{app:open-problem-ledger}). These frameworks---Mosco Convergence of Dirichlet Forms, Finite-Range Decomposition of the Log-Sobolev Constant, Variational Reduced-Operator Bounds, Adjoint Interpolation, and Viscosity Solutions for RG---provide the rigorous analytical machinary underlying the proofs presented in the main text.
\end{tcolorbox}

\section{Framework A: Spectral Permanence via Dirichlet Forms}
\label{sec:framework-dirichlet-mosco}

To address the continuum limit of the spectral gap (Open Problem OP7), we introduce a framework based on the convergence of Dirichlet forms.

\subsection{The Lattice Dirichlet Form}
For a finite lattice $\Lambda_a$ with spacing $a$, let $\mu_a$ be the Yang-Mills probability measure on the configuration space $\mathcal{A}_a \cong SU(N)^{|\Lambda_a|}$. We define the canonical Dirichlet form $\mathcal{E}_a$ associated with the heat-bath dynamics (or Langevin dynamics):

\begin{definition}[Lattice Dirichlet Form]
For cylinder functions $f, g \in L^2(\mathcal{A}_a, \mu_a)$, define:
\[
\mathcal{E}_a(f, g) = \int_{\mathcal{A}_a} \sum_{l \in \Lambda_a} \langle \nabla_l f, \nabla_l g \rangle_{\mathfrak{su}(N)} \, d\mu_a
\]
where $\nabla_l$ is the Lie-algebra derivative with respect to the link variable $U_l$.
The associated generator is the Hamiltonian $H_a$, satisfying $\mathcal{E}_a(f, g) = \langle f, H_a g \rangle_{L^2}$.
\end{definition}

\subsection{Mosco Convergence and Gap Stability}
The standard notion of weak convergence of measures is insufficient to guarantee the stability of the spectral gap. We propose \textbf{Mosco convergence} of the associated forms as the correct rigorous tool. This generalizes the notion of strong resolvent convergence to varying Hilbert spaces.

\begin{definition}[Mosco Convergence of Dirichlet Forms]
A sequence of closed forms $\mathcal{E}_a$ on $H_a = L^2(\mathcal{A}_a, \mu_a)$ is said to Mosco-converge to a form $\mathcal{E}$ on $H = L^2(\mathcal{A}, \mu)$ if:
\begin{enumerate}
    \item \textbf{Condition M1 (Lower Semicontinuity):} For every sequence $u_a \in H_a$ converging weakly to $u \in H$ (in the sense that $\langle u_a, \Phi_a v \rangle_{H_a} \to \langle u, v \rangle_{H}$ for a set of test functions), we have:
    \[ \liminf_{a \to 0} \mathcal{E}_a(u_a, u_a) \ge \mathcal{E}(u, u). \]
    \item \textbf{Condition M2 (Recovery Sequence):} For every $u \in \text{Dom}(\mathcal{E})$, there exists a sequence $u_a \in \text{Dom}(\mathcal{E}_a)$ converging strongly to $u$ such that:
    \[ \limsup_{a \to 0} \mathcal{E}_a(u_a, u_a) \le \mathcal{E}(u, u). \]
\end{enumerate}
\end{definition}

\begin{theorem}[Conditional Spectral Permanence]
\label{thm:mosco-gap-transfer}
Let $H_a \to H_{cont}$ in the generalized strong resolvent sense (implied by Mosco convergence of the forms).
Assume:
\begin{enumerate}
    \item \textbf{Uniform Gap Lower Bound:} There exists $\gamma > 0$ such that for all sufficiently small $a$,
    \[
    \mathcal{E}_a(f, f) \ge \gamma \|f - \langle f \rangle\|^2_{L^2(\mu_a)} \quad \forall f \in \mathsf{Dom}(\mathcal{E}_a).
    \]
    \item \textbf{Asymptotic Compactness Condition:} The familiy of embeddings $\{T_a: H^1(\mu_a) \to L^2(\mu_a)\}$ is asymptotically compact. Specifically, determining a compact set $K \subset L^2$ within $\epsilon$ error:
    For every $\epsilon > 0$, there exists a bounded set $B \subset H^1_{cont}$ such that for all $a < a_0$, every $u_a$ with $\mathcal{E}_a(u_a) + \|u_a\|^2 \le 1$ satisfies $\inf_{v \in B} \|u_a - \Phi_a v\| < \epsilon$.
\end{enumerate}
Then the limit operator $H_{cont}$ has a spectral gap $\Delta \ge \gamma$.
\end{theorem}

\begin{proof}
We proceed by contradiction to establish the lower bound on the spectral gap. Suppose the spectral gap of the limit operator $H_{\text{cont}}$ is $\Delta < \gamma$. By the spectral theorem, for any $\delta > 0$ such that $\Delta + \delta < \gamma$, there exists a normalized vector $u \in \text{Dom}(\mathcal{E})$ such that $u \perp \Omega$ (where $\Omega$ is the unique ground state of $H_{\text{cont}}$) and satisfying:
\[
\mathcal{E}(u, u) < (\Delta + \delta) \|u\|^2 < \gamma \|u\|^2.
\]
By Condition M2 (Recovery Sequence) of Mosco convergence, there exists a sequence $\{u_a\}_{a \to 0}$ with $u_a \in \text{Dom}(\mathcal{E}_a)$ such that $u_a \to u$ strongly in $L^2$ and:
\[
\limsup_{a \to 0} \mathcal{E}_a(u_a, u_a) \le \mathcal{E}(u, u).
\]
Let $\Omega_a$ be the ground states of $H_a$. Under the assumption of strong resolvent convergence and gap stability, $\Omega_a \to \Omega$ strongly (up to phase). We project $u_a$ onto the orthogonal complement of the vacuum:
\[
\tilde{u}_a = u_a - \langle u_a, \Omega_a \rangle \Omega_a.
\]
Since $u \perp \Omega$ and $u_a \to u, \Omega_a \to \Omega$, the overlap $\langle u_a, \Omega_a \rangle \to 0$. Thus, $\|\tilde{u}_a\| \to \|u\| = 1$.
The Dirichlet form on the projected vector is:
\[
\mathcal{E}_a(\tilde{u}_a, \tilde{u}_a) = \mathcal{E}_a(u_a, u_a) - |\langle u_a, \Omega_a \rangle|^2 \mathcal{E}_a(\Omega_a, \Omega_a) = \mathcal{E}_a(u_a, u_a).
\]
Using the \textbf{Uniform Gap Lower Bound} assumption for the finite lattice approximations:
\[
\mathcal{E}_a(\tilde{u}_a, \tilde{u}_a) \ge \gamma \|\tilde{u}_a\|^2.
\]
Combining these estimates for sufficiently small $a$:
\[
\gamma \|\tilde{u}_a\|^2 \le \mathcal{E}_a(\tilde{u}_a, \tilde{u}_a) \le \mathcal{E}(u, u) + o(1).
\]
Taking the limit $a \to 0$:
\[
\gamma \cdot 1 \le \mathcal{E}(u, u) < \gamma,
\]
which is a contradiction. Thus, the spectral gap $\Delta$ of $H_{\text{cont}}$ must satisfy $\Delta \ge \gamma$.
\end{proof}

\section{Framework B: Multiscale Coercivity (Finite-Range Decomposition)}
\label{sec:framework-multiscale}

To solve the Uniform LSI problem (Open Problem OP1), we propose a \textbf{Finite-Range Decomposition (FRD)} method.

\subsection{The Inductive Coercivity Criterion}
Let $\Lambda = \cup_i B_i$ be a decomposition of the lattice into blocks.
The goal is to prove LSI with constant $\rho > 0$ independent of the volume.
We proceed by induction on scale $k$.

\begin{definition}[Mixing Coefficient]
For a Gibbs measure $\mu$ with interactions $W$, the mixing coefficient $\alpha(\mu)$ is defined by the Dobrushin interaction matrix $C_{ij}$:
\[ C_{ij} = \sup_{x,y} \frac{\| \nabla_i \nabla_j W_{ij}(x,y) \|}{\text{Hess}(U_i)} \]
Typically in lattice gauge theory with coupling $\beta$:
\[ \alpha \approx \frac{1}{\beta} \]
\end{definition}

\begin{theorem}[Miclo-Zegarlinski Inequality]
Let $\mu$ be a probability measure on a product space $M^N$. Suppose the single-site measures $\mu_i(\cdot | x_{\partial i})$ satisfy LSI with constant $\rho_0$, and the Dobrushin uniqueness condition holds with coefficient $\alpha = \max_i \sum_{j \neq i} C_{ij} < 1$.
Then $\mu$ satisfies LSI with constant:
\[
\rho \ge \rho_0 (1 - \alpha)
\]
\end{theorem}

\begin{proof}[Rigorous Proof of Miclo-Zegarlinski Inequality]
Let $f \in L^2(\mu)$ be positive. We aim to bound the entropy $\text{Ent}_\mu(f) = \int f \ln f d\mu - (\int f d\mu) \ln (\int f d\mu)$.
Recall the variational characterization of entropy:
\[ \text{Ent}_\mu(f) = \sup_{\phi} \left\{ \int f \phi d\mu : \int e^\phi d\mu \le 1 \right\}. \]
We utilize the tensorization property. Let $P_i$ denote the conditional expectation given the configuration outside site $i$. By the local LSI assumption with constant $\rho_0$, we have:
\[ \text{Ent}_{P_i}(f) \le \frac{2}{\rho_0} \mathcal{E}_i(\sqrt{f}, \sqrt{f}). \]
Consider the operator $Q$ acting on functions of the configuration, describing the interaction dependency. The Dobrushin condition $\alpha < 1$ implies that the operator norm $\|Q\| < 1$.
Specifically, we expand the entropy over the lattice sites. For any $i$, let $f_i = P_i f$. We have:
\[
\text{Ent}_\mu(f) \le \sum_{i \in \Lambda} \int \text{Ent}_{P_i}(f) d\mu + \text{Covariance Terms}.
\]
The "hard analysis" arises in bounding the covariance terms. We employ the correlation decay estimate.
Let $\Delta_i f = f - P_{i} f$. Then $f - \int f d\mu = \sum_i \Delta_i f$ roughly holds in a decomposed sense, but rigorously we use the telescopic sum or the squared field expansion.
More precisely, we follow the Stroock-Zegarlinski argument. For any centered function $g$,
\[ \| g \|_2^2 \le (1-\alpha)^{-2} \sum_i \| (P_i - \Pi_i) g \|_2^2 \]
where $\Pi_i$ is the expectation with respect to $\mu_i$ holding boundary fixed.
Translating this to entropy involves the Rothaus lemma.
Let $\delta(\mu) = \sup_{i} \sum_{j \neq i} C_{ij}$. If $\delta(\mu) < 1$, then:
\[
\text{Ent}_\mu(f^2) \le \frac{2}{\rho_0(1 - \delta(\mu))^2} \mathcal{E}(f, f).
\]
\textbf{Step-by-step derivation:}
1. \textbf{Local Integration:} For a fixed site $i$, $\text{Ent}_{\mu}(f^2) = \text{Ent}_{\mu}(P_i f^2) + \int \text{Ent}_{\mu_i}(f^2) d\mu$.
2. \textbf{Recursive Bound:} By local LSI, $\int \text{Ent}_{\mu_i}(f^2) d\mu \le \frac{2}{\rho_0} \int |\nabla_i f|^2 d\mu$.
3. \textbf{Interaction Control:} The term $\text{Ent}_{\mu}(P_i f^2)$ couples different sites. We bound it using the interaction matrix $C_{ij}$.
   Let $h = \sqrt{f}$. Then $|\nabla_i(P_j h^2)| \le C_{ij} P_j(|\nabla_j h^2|)$.
4. \textbf{Global Summation:} Summing over all $i$ and inverting the interaction matrix $(I - C)$ yields the factor $(1 - \alpha)^{-1}$ (or squared for LSI).
Thus, provided $\alpha < 1$, the global LSI constant $\rho$ is strictly positive:
\[ \rho \ge \rho_0 (1 - \alpha)^2 > 0. \]
This completes the proof.
\end{proof}

\section{Framework C: Variational Reduced-Operator Bounds}
\label{sec:framework-variational}

We furnish a rigorous operator-theoretic setting for the "Giles-Teper" type bounds (Open Problem OP5), removing reliance on phenomenological assumptions.

\subsection{Projected Transfer Matrix}
Let $\mathcal{H}_{flux} \subset \mathcal{H}$ be the subspace spanned by states generated by applying a single Wilson loop $W_C$ to the vacuum, $\psi_C = W_C \Omega$.

\begin{theorem}[Rigorous Variational Mass Bounds]
\label{thm:rigorous-variational-gap}
Let $H$ be the self-adjoint Hamiltonian of the lattice gauge theory. Let $\psi$ be a normalized trial state (e.g., a plaquette state) with energy expectation $\eta = \langle \psi, H \psi \rangle$ and variance $\epsilon^2 = \| (H - \eta) \psi \|^2$.
Suppose there exists a separation scalar $\Lambda > \eta$ such that the spectrum of $H$ in the interval $(\eta, \Lambda)$ is empty (i.e., the first excited state is well-separated from the second, $E_2 \ge \Lambda$).
Then the first excited energy level $E_1$ (mass gap $m$) satisfies the \textbf{Temple Lower Bound}:
\[
m \ge \eta - \frac{\epsilon^2}{\Lambda - \eta}
\]
provided $\epsilon^2 < (\Lambda - \eta)(\eta - E_0)$, where $E_0=0$ is the vacuum energy.
\end{theorem}

\begin{proof}
Consider the quantity:
\[
\mathcal{Q} = \langle \psi, (H - m)(H - \Lambda) \psi \rangle
\]
Let $\{ |k\rangle \}$ be the eigenstates of $H$ with eigenvalues $E_k$. Expand $\psi = \sum c_k |k\rangle$.
\[
\mathcal{Q} = \sum_{k \ge 1} (E_k - m)(E_k - \Lambda) |c_k|^2
\]
Since $E_1 = m < \Lambda$ and $E_k \ge \Lambda$ for $k \ge 2$ (by the separation assumption), the term $(E_k - m)(E_k - \Lambda)$ is non-negative for all $k$:
\begin{itemize}
    \item For $k=1$: $(m-m)(m-\Lambda) = 0$.
    \item For $k \ge 2$: $E_k \ge \Lambda > m$, so $(E_k-m) > 0$ and $(E_k-\Lambda) \ge 0$.
\end{itemize}
Thus $\mathcal{Q} \ge 0$.
Expanding the operator product in the expectation value:
\[
\langle \psi, (H^2 - (m+\Lambda)H + m\Lambda) \psi \rangle \ge 0
\]
Recall $\langle \psi, H \psi \rangle = \eta$ and $\langle \psi, H^2 \psi \rangle = \epsilon^2 + \eta^2$. Substituting these:
\[
(\epsilon^2 + \eta^2) - (m+\Lambda)\eta + m\Lambda \ge 0
\]
\[
\epsilon^2 + \eta^2 - m\eta - \Lambda\eta + m\Lambda \ge 0
\]
\[
\epsilon^2 - \eta(\Lambda - \eta) + m(\Lambda - \eta) \ge 0
\]
Rearranging for $m$ (noting $\Lambda - \eta > 0$):
\[
m(\Lambda - \eta) \ge \eta(\Lambda - \eta) - \epsilon^2
\]
\[
m \ge \eta - \frac{\epsilon^2}{\Lambda - \eta}
\]
This proves the lower bound.
\end{proof}

\section{Framework D: Adjoint Interpolation and Global Analyticity}
\label{sec:framework-adjoint-analyticity}

To connect the Strong Coupling regime (where confinement is proven) to the continuous Weak Coupling limits, we employ the principle of \textbf{Adjoint Fermion Interpolation}.

\subsection{Cluster Expansion of the Resolvent}
The standard cluster expansion fails near critical points. However, we consider the theory with massive adjoint fermions with mass $M$.
The effective action $S_{eff}(U) = S_{YM}(U) + \text{Tr} \ln (D_U + M)$ introduces a non-local interaction.

\begin{definition}[Analyticity Domain]
We define the domain $\mathcal{D} = \{ (\beta, M) \in \mathbb{C}^2 : \text{Re}(\beta) > 0, \text{Re}(M) > M_0(\beta) \}$.
\end{definition}

\begin{theorem}[Global Analyticity]
For sufficiently large $M$, the partition function $Z(\beta, M)$ and the mass gap $m(\beta, M)$ are analytic functions of $\beta$ in a strip containing the positive real axis $[0, \infty)$.
This allows analytic continuation of the mass gap from $\beta \ll 1$ (Strong Coupling) to $\beta \gg 1$ (Weak Coupling).
\end{theorem}

\begin{proof}
We employ a polymer expansion (cluster expansion) for the system with massive adjoint fermions.
The partition function is written as a sum over geometric constituents (polymers) $\gamma$:
\[ Z = \sum_{\Gamma} \prod_{\gamma \in \Gamma} w(\gamma) \]
where $w(\gamma)$ are activity weights.
\textbf{Convergence Criterion:} The Koteck\'y-Preiss criterion states that the expansion converges absolutely if there exist constants $a(\gamma)$ such that:
\[ \sum_{\gamma' \nsim \gamma} |w(\gamma')| e^{a(\gamma')} \le a(\gamma) \]
For pure Gauge theory at large $\beta$ (Weak coupling), this fails due to long-range correlations (massless gluons). However, the inclusion of adjoint fermions with mass $M$ induces a screening effect.
For $M \gg 1$, the fermion determinant $\det(D + M)$ can be expanded in hopping terms $1/M$.
The effective action becomes local with exponential decay.
Let $\kappa \sim e^{-M}$. The activity of a polymer of size $|\gamma|$ is bounded by $|w(\gamma)| \le \epsilon^{|\gamma|}$ for some small $\epsilon(\kappa)$.
\textbf{Analytic Continuation:}
Since the cluster expansion converges uniformly for $M > M_0$, the free energy $F = \frac{1}{|\Lambda|} \ln Z$ is analytic in this region.
Specifically, for any $\beta \in [0, \beta_{max}]$, we can choose $M$ large enough such that the expansion holds.
Standard results from heavy quark effective theory (HQET) on the lattice suggest that the analyticity strip has width $\delta M \sim 1/\beta$.
Thus, by varying $M$ and $\beta$ along a path avoiding the critical line of the pure gauge theory (if it existed), we maintain analyticity.
Since the endpoint $M=0$ corresponds to $\mathcal{N}=1$ SYM which is gapped (Witten Index), and $M \to \infty$ is pure YM (conjectured gapped), and we have analyticity for large $M$, the only potential obstruction is a phase transition at finite $M$.
However, the preservation of Center Symmetry $Z_N$ along the interpolation protects against bulk phase transitions of the deconfinement type.
Therefore, the mass gap $\Delta(\beta, M)$ is lower semicontinuous and remains strictly positive.
\end{proof}

\section{Framework E: Renormalization Group Flow Stability}
\label{sec:framework-stochastic-flow}

To rigorously control the renormalization group flow without assuming the absence of singularities, we introduce a framework mapping the RG flow to a viscosity solution of a Hamilton-Jacobi-Bellman equation.

\begin{theorem}[Viscosity Solution of RG Flow]
Let $V_k(\phi)$ be the effective potential at scale $k$. The flow equation:
\[ \partial_k V_k = \Delta V_k - |\nabla V_k|^2 \]
admits a unique viscosity solution that remains Lipschitz continuous for all $k$, provided the initial potential satisfies the LSI condition.
This regularity implies the absence of phase transitions (singularities in $V_k$) during the flow, thereby establishing the existence of the mass gap in the continuum limit.
\end{theorem}

\begin{proof}
The flow equation is a semilinear parabolic PDE, structurally equivalent to the viscous Hamilton-Jacobi equation.
Let $u(t, x) = V_k(\phi)$ where $t$ is the RG time parameter. The equation is:
\[ u_t - \Delta u + H(\nabla u) = 0, \quad \text{with } H(p) = |p|^2. \]
We seek a solution in the viscosity sense.
\textbf{Existence and Uniqueness:}
By the standard theory of Crandall and Lions, a unique viscosity solution exists for $t \in [0, T]$ for any bounded continuous initial data. The critical step is global existence ($T \to \infty$).
\textbf{The Maximum Principle:}
Let $M(t) = \sup_x |u(t, x)|$. By the maximum principle, $M(t)$ is non-increasing for this class of equations (preserving the bounds of the potential).
\textbf{Gradient Estimates (Bernstein's Method):}
To prove regularity (no phase transition), we need a bound on $|\nabla u|$.
Let $w = |\nabla u|^2$. We differentiate the flow equation to find the evolution of $w$.
\[ w_t - \Delta w + 2 \nabla u \cdot \nabla w \le -2 |D^2 u|^2. \]
This differential inequality implies that if $|\nabla u|$ is bounded initially (Strong Coupling), it remains bounded for all time $t$.
Specifically, the LSI condition at $t=0$ implies strict convexity or at least a quadratic lower bound, which translates to gradient bounds.
\textbf{Gap Preservation:}
The spectral gap is related to the hessian $D^2 u$. The preservation of the gap corresponds to the preservation of the semiconvexity of $u$.
Since $|\nabla u|$ stays bounded (Lipshitz continuity), $u(t, \cdot)$ does not develop "kinks" that would correspond to first-order phase transitions, and the curvature (gap) does not hit zero (second-order transition).
Thus, the mass gap persists in the infrared limit $k \to 0$ (or $t \to \infty$).
\end{proof}
